%%% chapters/02_introduction.tex — 서론
\section{서론}\label{sec:introduction}

\subsection{연구 배경}\label{ssec:background}

프로그래밍 언어 간의 거리를 정량적으로 측정하려는 시도는 프로그래밍 언어론과 소프트웨어 공학의 오랜 주제이다. 
전통적으로 언어 간 유사도는 구문(syntax) 계열, 타입 시스템, 패러다임(절차적, 객체지향적, 함수형) 등의 정성적 분류에 의존해 왔다.
그러나 이러한 분류는 실제 코드 변환 과정에서 나타나는 구조적, 의미적 변형의 정도를 정확히 반영하지 못한다.

한편, 대규모 언어 모델(LLM)의 코드 생성 및 번역 능력이 급속히 발전하면서, LLM을 번역기로 활용하여 언어 간 변환을 수행하고 그 과정에서 발생하는 변형을 측정하는 새로운 접근이 가능해졌다.
자연어 처리 분야에서 왕복 번역(round-trip translation)은 기계 번역 품질을 평가하는 고전적 기법이며, 이를 프로그래밍 언어에 적용하면 언어 간 거리를 반영하는 정량적 지표를 도출할 수 있다는 가설이 본 연구의 출발점이다.

\subsection{연구 목적}\label{ssec:objectives}

본 연구의 목적은 다음 세 가지로 요약된다.

\begin{enumerate}[leftmargin=2em, label=\arabic*)]
  \item \textbf{RTT 기반 언어 거리 측정 프레임워크 구축}:
    \cpp{}를 기준 언어(seed language)로 설정하고, C\, Java, Python으로 
    왕복 번역 순환을 통해 언어 간 거리를 다차원적으로 측정하는 자동화된 실험
    파이프라인을 설계, 구현한다.
  \item \textbf{이종 LLM 간 비교 분석}:
    로컬 오픈소스 모델(qwen2.5-coder:7b)과 상용 클라우드 모델(gpt-5.4)의 번역 행동을 동일 코퍼스\, 동일 조건에서 비교하여, 모델 특성이 언어 간 거리 측정 결과에 미치는 영향을 분석한다.
  \item \textbf{다차원 지표의 종합적 해석}:
  	수렴 속도(반복 횟수), 잔차 토큰 유사도, AST 편집 거리, MOSS 표절 검사 유사도,
    의미 보존 여부(semantic pass/fail) 등 이질적 지표들의 상호 관계를 분석하여,
    각 지표가 드러내는 언어 간 거리의 서로 다른 측면을 조명한다.
\end{enumerate}